\chapter{Maximum Principle and Li-Yau-Hamilton Inequalities}
\label{chap:cz-maximum-principles}

\section{Preserving Positive Curvature}

Throughout this section, $g(t)$ is a complete Ricci flow on $[0,T]$
with uniformly bounded curvature.  The curvature operator evolves by
\[
\partial_t M_{\alpha\beta}=\Delta M_{\alpha\beta}
+M_{\alpha\gamma}M_{\beta\gamma}+M_{\alpha\beta}^{\#}.
\]

\begin{lemma}[proper exhaustion with uniformly bounded derivatives]
  \label{lem:cz-proper-exhaustion-with-uniformly-bounded-derivatives}
  \source{cao-zhu}
  \dcref{Lemma 2.1.1}
  \group{Cao--Zhu Preserving Positive Curvature}
There exists a smooth function $f$ on $M$ such that $f\ge 1$
everywhere, $f(x)\rightarrow +\infty$ as $d_0(x,x_0)\rightarrow
+\infty$ $($for some fixed $x_0\in M),$
$$
|\nabla f|_{g_{ij}(t)}\le C \quad \text{and}\quad |\nabla^2
f|_{g_{ij}(t)}\le C
$$
on $M \times [0,T]$ for some positive constant $C$.
\end{lemma}

\begin{proof}
  \uses{thm:cz-global-curvature-derivative-estimates}
  \sketch Smooth a distance exhaustion at time zero, use Hessian comparison, and propagate its first two derivative bounds with the Ricci-flow connection and curvature derivative estimates.
\end{proof}

\begin{proposition}[preservation of nonnegative scalar curvature]
  \label{prop:cz-preservation-of-nonnegative-scalar-curvature}
  \source{cao-zhu}
  \dcref{Proposition 2.1.2}
  \group{Cao--Zhu Preserving Positive Curvature}
If the scalar curvature $R$ of the solution $g_{ij}(t), 0\le t\le
T$, to the Ricci flow is nonnegative at $t=0$, then it remains so
on $0\le t\le T$.
\end{proposition}

\begin{proof}
  \uses{lem:cz-proper-exhaustion-with-uniformly-bounded-derivatives}
  \sketch Apply the scalar maximum principle to $R+\varepsilon e^{At}f$ and let $\varepsilon\to0$.
\end{proof}

\begin{lemma}[tensor maximum principle on complete manifolds]
  \label{lem:cz-tensor-maximum-principle-on-complete-manifolds}
  \source{cao-zhu}
  \dcref{Lemma 2.1.3}
  \group{Cao--Zhu Preserving Positive Curvature}
Suppose that on $0\le t\le T$,
$$
\frac{\partial}{\partial t}M_{\alpha\beta} =\Delta
M_{\alpha\beta}+u^i\nabla_i M_{\alpha\beta}+N_{\alpha\beta}
$$
where $u^i(t)$ is a time-dependent vector field on $M$ with
uniform bound and
$N_{\alpha\beta}={\mathcal{P}}(M_{\alpha\beta},h_{\alpha\beta})$
satisfies
$$
N_{\alpha\beta}v^\alpha v^\beta\ge 0 \quad \text{whenever} \quad
M_{\alpha\beta}v^\beta=0.
$$
If $M_{\alpha\beta}\ge 0$ at $t=0$, then it remains so on $0\le
t\le T$.
\end{lemma}

\begin{proof}
  \uses{lem:cz-proper-exhaustion-with-uniformly-bounded-derivatives}
  \sketch Perturb $M$ by $\varepsilon e^{At}fh$; a first null vector contradicts the null-eigenvector condition when $A$ is large.
\end{proof}

\begin{proposition}[preservation of nonnegative curvature operator]
  \label{prop:cz-preservation-of-nonnegative-curvature-operator}
  \source{cao-zhu}
  \dcref{Proposition 2.1.4}
  \group{Cao--Zhu Preserving Positive Curvature}
Nonnegativity of the curvature operator $M_{\alpha\beta}$ is
preserved by the Ricci flow.
\end{proposition}

\begin{proposition}[preservation of holomorphic bisectional curvature]
  \label{prop:cz-preservation-of-holomorphic-bisectional-curvature}
  \source{cao-zhu}
  \dcref{Proposition 2.1.5}
  \group{Cao--Zhu Preserving Positive Curvature}
Under the K\"ahler-Ricci flow if the initial metric has positive
$($nonnegative$)$ holomorphic bisectional curvature then the
evolved metric also has positive $($nonnegative$)$ holomorphic
bisectional curvature.
\end{proposition}

\section{Strong Maximum Principle}

Let $V\to\Omega$ have a fixed bundle metric and a metric-compatible,
possibly time-dependent connection.  For a symmetric form $M$ consider
\[
\partial_tM=\Delta M+N(M), \qquad N(M)\ge0\ \text{when}\ M\ge0.
\tag{2.2.1}
\]

\begin{theorem}[strong maximum principle for symmetric bilinear forms]
  \label{thm:cz-strong-maximum-principle-for-symmetric-bilinear-forms}
  \source{cao-zhu}
  \dcref{Theorem 2.2.1}
  \group{Cao--Zhu Strong Maximum Principle}
 Let
$M_{\alpha\beta}$ be a smooth solution of the equation $(2.2.1)$.
Suppose $M_{\alpha\beta}\geq0$ on $\Omega \times [0,T]$. Then
there exists a positive constant $0<\delta \leq T$ such that on
$\Omega \times (0,\delta)$, the rank of $M_{\alpha\beta}$ is
constant, and the null space of $M_{\alpha\beta}$ is invariant
under parallel translation and invariant in time and also lies in
the null space of $N_{\alpha\beta}$.
\end{theorem}

\begin{proof}
  \sketch Use the heat evolution of a rank-detecting barrier to obtain constant rank, then differentiate $M(v,\cdot)=0$ to prove the asserted null-space invariances.
\end{proof}

\begin{theorem}[parallel invariance of the curvature-operator image]
  \label{thm:cz-parallel-invariance-of-the-curvature-operator-image}
  \source{cao-zhu}
  \dcref{Theorem 2.2.2}
  \group{Cao--Zhu Strong Maximum Principle}
Suppose the curvature operator $M_{\alpha\beta}$ of the initial
metric is nonnegative. Then, under the Ricci flow, for some
interval $0<t<\delta$ the image of $M_{\alpha\beta}$ is a Lie
subalgebra of $so(n)$ which has constant rank and is invariant
under parallel translation and invariant in time.
%$$\eqno\#$$ \vskip 1cm
\end{theorem}

\section{Advanced Maximum Principle for Tensors}

Let $V\to M$ have a fixed bundle metric, time-dependent compatible
connections $\nabla_t$, and a bounded-drift nonlinear equation
\[
\partial_t\sigma=\Delta_t\sigma+u^i(\nabla_t)_i\sigma+N(x,\sigma,t).
\tag{PDE}
\]
The associated fiber equation is
\[
\dot\sigma_x=N(x,\sigma_x,t).
\tag{ODE}
\]
For a closed set $K\subset V$, (H1) means parallel-translation
invariance and (H2) means that every $K_x$ is closed and convex.

\begin{theorem}[advanced maximum principle for time-dependent tensors]
  \label{thm:cz-advanced-maximum-principle-for-time-dependent-tensors}
  \source{cao-zhu}
  \dcref{Theorem 2.3.1}
  \group{Cao--Zhu Advanced Maximum Principle for Tensors}
 Let $K$ be a closed subset of $V$ satisfying the
hypothesis {\rm (H1)} and {\rm (H2)}. Suppose that for any $x\in
M$ and any initial time $t_0\in[0,T)$, any solution $\sigma_x(t)$
of the {\rm (ODE)} which starts in $K_x$ at $t_0$ will remain in
$K_x$ for all later times. Then for any initial time $t_0\in
[0,T)$ the solution $\sigma(x,t)$ of the {\rm (PDE)} will remain
in $K$ for all later times provided $\sigma(x,t)$ starts in $K$ at
time $t_0$ and $\sigma(x,t)$ is uniformly bounded with respect to
the bundle metric $h_{ab}$ on $M\times[t_0,T]$.
\end{theorem}

\begin{lemma}[vanishing from a Lipschitz differential inequality]
  \label{lem:cz-vanishing-from-a-lipschitz-differential-inequality}
  \source{cao-zhu}
  \dcref{Lemma 2.3.2}
  \group{Cao--Zhu Advanced Maximum Principle for Tensors}
Suppose $\varphi: [a,b]\rightarrow \mathbb{R}$ is Lipschitz
continuous and suppose for some constant $C<+\infty$,
\begin{align*}
&  \frac{d}{dt}\varphi(t)\le C\varphi(t),\quad whenever\ \
\varphi(t)\ge 0\;\ on\ \;[a,b),\\[2mm]
\text{and}\hskip 1cm &  \varphi(a)\le 0.
\end{align*}
Then $\varphi(t)\le 0$ on $[a,b]$.
\end{lemma}

\begin{proof}
  \sketch After multiplying by $e^{-Ct}$, compare with $\varepsilon(t-a)$ and let $\varepsilon\to0$.
\end{proof}

\begin{lemma}[derivative of a fiberwise supremum]
  \label{lem:cz-derivative-of-a-fiberwise-supremum}
  \source{cao-zhu}
  \dcref{Lemma 2.3.3}
  \group{Cao--Zhu Advanced Maximum Principle for Tensors}
$$
\frac{d}{dt}\varphi(t) \le \sup\left\{\frac{\partial
\psi}{\partial t}(y,t)\ |\ y\in Y \; \mbox{ satisfies }\;
\psi(y,t)=\varphi(t)\right\}.
$$
\end{lemma}

\begin{proof}
  \sketch Choose maximizing points along a sequence of forward times, pass to a compact-limit maximizer, and apply the mean value theorem.
\end{proof}

\begin{lemma}[invariance criterion for closed convex sets]
  \label{lem:cz-invariance-criterion-for-closed-convex-sets}
  \source{cao-zhu}
  \dcref{Lemma 2.3.4}
  \group{Cao--Zhu Advanced Maximum Principle for Tensors}
Let $U\subset R^n$ be an open set and $Z\subset U$ be a closed
convex subset. Consider the {\rm ODE} \be
\frac{d\varphi}{dt}=N(\varphi,t)%\eqno(2.3.1)
\ee where $N:\ U\times[0,T]\rightarrow R^n$ is continuous and
Lipschitz in $\varphi$. Then the following two statements are
equivalent.
\begin{itemize}
\item[(i)] For any initial time $t_0\in[0,T]$, any solution of the
{\rm ODE} $(2.3.1)$ which starts in Z at $t_0$ will remain in Z
for all later times; \item[(ii)] $\varphi+N(\varphi,t)\in
T_\varphi Z$ for all $\varphi\in\partial Z$ and $t\in[0,T)$.
\end{itemize}
\end{lemma}

\begin{proof}
  \uses{lem:cz-derivative-of-a-fiberwise-supremum,lem:cz-vanishing-from-a-lipschitz-differential-inequality}
  \sketch Use support functionals to identify the tangent-cone condition; the distance to $Z$ then obeys the preceding Lipschitz inequality.
\end{proof}

\begin{theorem}[maximum principle for closed fiberwise convex sets]
  \label{thm:cz-maximum-principle-for-closed-fiberwise-convex-sets}
  \source{cao-zhu}
  \dcref{Theorem 2.3.5}
  \group{Cao--Zhu Advanced Maximum Principle for Tensors}
Let $K(t)\subset V$, $t\in [0,T]$ be closed subsets which satisfy
the following hypotheses
\begin{itemize}
\item[(H3)] $K(t)$ is invariant under parallel translation defined
by the connection $\nabla_t$ for each $t\in [0,T]$; \item[(H4)] in
each fiber $V_x$, the set $K_x(t)\stackrel{\Delta}{=}K(t)\cap V_x$
is nonempty, closed and convex for each $t\in [0,T]$; \item[(H5)]
the space-time track $\bigcup\limits_{t\in[0,T]}(\partial
K(t)\times\{t\})$ is a closed subset of $V\times[0,T]$.
\end{itemize}
Suppose that, for any $x\in M$ and any initial time $t_0\in[0,T)$,
and for any solution $\sigma_x(t)$ of the {\rm (ODE)} which starts
in $K_x(t_0)$, the solution $\sigma_x(t)$ will remain in $K_x(t)$
for all later times. Then for any initial time $t_0\in[0,T)$ the
solution $\sigma(x,t)$ of the {\rm (PDE)} will remain in $K(t)$
for all later times if $\sigma(x,t)$ starts in $K(t_0)$ at time
$t_0$ and the solution $\sigma(x,t)$ is uniformly bounded with
respect to the bundle metric $h_{ab}$ on $M\times[t_0,T]$.
\end{theorem}

\section{Hamilton-Ivey Curvature Pinching Estimate}

In dimension three, write the curvature-operator eigenvalues as
$\lambda\ge\mu\ge\nu$; their ODE is
\[
\dot\lambda=\lambda^2+\mu\nu,\quad
\dot\mu=\mu^2+\lambda\nu,\quad
\dot\nu=\nu^2+\lambda\mu.
\]

\begin{theorem}[Hamilton–Ivey curvature pinching]
  \label{thm:cz-hamiltonivey-curvature-pinching}
  \source{cao-zhu}
  \dcref{Theorem 2.4.1}
  \group{Cao--Zhu Hamilton-Ivey Curvature Pinching Estimate}
Suppose we have a solution to the Ricci flow on a three-manifold
which is complete with bounded curvature for each $t\geq0$. Assume
at $t=0$ the eigenvalues $\lambda\geq\mu\geq\nu$ of the curvature
operator at each point are bounded below by $\nu\geq-1$. The
scalar curvature $R=\lambda+\mu+\nu$ is their sum. Then at all
points and all times $t\geq0$ we have the pinching estimate
$$
R\geq (-\nu)[\log(-\nu)-3],
$$
whenever $\nu<0$.
\end{theorem}

\begin{proof}
  \uses{thm:cz-advanced-maximum-principle-for-time-dependent-tensors}
  \sketch The closed convex region $R\ge-3$ and $\nu+f^{-1}(R)\ge0$ is preserved by the curvature-operator ODE, hence by the advanced maximum principle.
\end{proof}

\section{Li-Yau-Hamilton Estimates}

For a Ricci flow, set
\[
P_{ijk}=\nabla_iR_{jk}-\nabla_jR_{ik},\qquad
M_{ij}=\Delta R_{ij}-\tfrac12\nabla_i\nabla_jR
+2R_{ikjl}R_{kl}-R_{ik}R_{jk}+\tfrac1{2t}R_{ij}.
\]

\begin{theorem}[Li–Yau differential Harnack estimate]
  \label{thm:cz-liyau-differential-harnack-estimate}
  \source{cao-zhu}
  \dcref{Theorem 2.5.1}
  \group{Cao--Zhu Li-Yau-Hamilton Estimates}
Let $(M, g_{ij})$ be an $n$-dimensional complete Riemannian
manifold with nonnegative Ricci curvature. Let $u(x,t)$ be any
positive solution to the heat equation
$$
\frac{\partial u}{\partial t}=\Delta u \qquad {\text{on}} \ \
M\times [0,\infty).
$$
Then we have \be \frac{\partial u}{\partial t}-\frac{|\nabla
u|^2}{u}+\frac{n}{2t}u\ge 0 \qquad \text{on } \ \ M\times
(0,\infty).%\eqno(2.5.1)
\ee
%$$\eqno \#$$ \vskip 0.3cm
\end{theorem}

\begin{theorem}[surface Ricci-flow differential Harnack estimate]
  \label{thm:cz-surface-ricci-flow-differential-harnack-estimate}
  \source{cao-zhu}
  \dcref{Theorem 2.5.2}
  \group{Cao--Zhu Li-Yau-Hamilton Estimates}
Let $g_{ij}(x,t)$ be a complete solution of the Ricci flow on a
surface $M$. Assume the scalar curvature of the initial metric is
bounded, nonnegative everywhere and positive somewhere. Then the
scalar curvature $R(x,t)$ satisfies the Li-Yau estimate \be
\frac{\partial R}{\partial t} - \frac{|\nabla R|^2}{R} +
\frac{R}{t} \geq 0.%\eqno(2.5.4)
\ee
\end{theorem}

\begin{proof}
  \uses{lem:cz-tensor-maximum-principle-on-complete-manifolds}
  \sketch Apply the tensor maximum principle to the surface Harnack quantity using the scalar-curvature evolution.
\end{proof}

\begin{corollary}[surface Ricci-flow integrated Harnack estimate]
  \label{cor:cz-surface-ricci-flow-integrated-harnack-estimate}
  \source{cao-zhu}
  \dcref{Corollary 2.5.3}
  \group{Cao--Zhu Li-Yau-Hamilton Estimates}
Let $g_{ij}(x,t)$ be a complete solution of the Ricci flow on a
surface with bounded and nonnegative scalar curvature. Then for
any points $x_{1},x_{2}\in M$, and $0<t_{1}<t_{2}$, we have
$$
R(x_2,t_2)
\geq\frac{t_1}{t_2}e^{-d_{t_1}{(x_1,x_2)}^2/{4(t_2-t_1)}}R(x_1,t_1).
$$
\end{corollary}

\begin{proof}
  \uses{thm:cz-surface-ricci-flow-differential-harnack-estimate}
  \sketch Integrate the differential Harnack inequality along a minimizing $g(t_1)$-geodesic joining $x_1$ to $x_2$.
\end{proof}

\begin{theorem}[matrix Li–Yau–Hamilton estimate]
  \label{thm:cz-matrix-liyauhamilton-estimate}
  \source{cao-zhu}
  \dcref{Theorem 2.5.4}
  \group{Cao--Zhu Li-Yau-Hamilton Estimates}
Let $g_{ij}(x,t)$ be a complete solution with bounded curvature to
the Ricci flow on a manifold $M$ for $t$ in some time interval
$(0,T)$ and suppose the curvature operator of $g_{ij}(x,t)$ is
nonnegative. Then for any one-form $W_a$ and any two-form $U_{ab}$
we have
$$
M_{ab}W_aW_b+2P_{abc}U_{ab}W_c+R_{abcd}U_{ab}U_{cd}\geq 0
$$
on $M\times (0,T)$.
\end{theorem}

\begin{corollary}[trace Li–Yau–Hamilton estimate]
  \label{cor:cz-trace-liyauhamilton-estimate}
  \source{cao-zhu}
  \dcref{Corollary 2.5.5}
  \group{Cao--Zhu Li-Yau-Hamilton Estimates}
For any one-form $V_a$ we have
$$
\frac{\partial R}{\partial t}+\frac{R}{t} +2\nabla_aR\cdot
V_a+2R_{ab}V_aV_b\geq 0.
$$
\end{corollary}

\begin{corollary}[rigidity in the trace Harnack estimate]
  \label{cor:cz-rigidity-in-the-trace-harnack-estimate}
  \source{cao-zhu}
  \dcref{Corollary 2.5.6}
  \group{Cao--Zhu Li-Yau-Hamilton Estimates}
Suppose we have a solution to the Ricci flow for $t>0$ which is
complete with bounded curvature, and has nonnegative curvature
operator. Suppose also that at some time $t>0$ we have the scalar
curvature $R\leq M$ for some constant M in the ball of radius $r$
around some point $p$. Then for $k=1,2,\ldots$, the $k^{th}$ order
derivatives of the curvature at $p$ at the time $t$ satisfy a
bound
$$
|\nabla^k Rm(p,t)|^2\leq
C_kM^2\left(\frac{1}{r^{2k}}+\frac{1}{t^k}+M^k\right)
$$
for some constant $C_k$ depending only on the dimension and $k$.
\end{corollary}

\begin{proof}
  \uses{thm:cz-local-curvature-derivative-estimates}
  \sketch Invoke the stated local curvature derivative estimate at the prescribed scale.
\end{proof}

\begin{corollary}[integrated scalar-curvature Harnack estimate]
  \label{cor:cz-integrated-scalar-curvature-harnack-estimate}
  \source{cao-zhu}
  \dcref{Corollary 2.5.7}
  \group{Cao--Zhu Li-Yau-Hamilton Estimates}
Let $g_{ij}(x,t)$ be a complete solution of the Ricci flow on a
manifold with bounded and nonnegative curvature operator, and let
$x_1,x_2\in M,\ 0<t_1<t_2.$ Then the following inequality holds
$$
R(x_2,t_2)\geq\frac{t_1}{t_2}e^{-d_{t_1}(x_1,x_2)^2/2(t_2-t_1)}\cdot
R(x_1,t_1).$$
\end{corollary}

\begin{corollary}[scalar curvature monotonicity for ancient solutions]
  \label{cor:cz-scalar-curvature-monotonicity-for-ancient-solutions}
  \source{cao-zhu}
  \dcref{Corollary 2.5.8}
  \group{Cao--Zhu Li-Yau-Hamilton Estimates}
Let $g_{ij}(x,t)$ be a complete ancient solution of the Ricci flow
on $M\times (-\infty, T)$ with bounded and nonnegative curvature
operator, then the scalar curvature $R(x,t)$ is pointwise
nondecreasing in time $t$.
\end{corollary}

\begin{theorem}[Kähler–Ricci matrix Harnack estimate]
  \label{thm:cz-kahlerricci-matrix-harnack-estimate}
  \source{cao-zhu}
  \dcref{Theorem 2.5.9}
  \group{Cao--Zhu Li-Yau-Hamilton Estimates}
Let $g_{\alpha\bar \beta}(x,t)$ be a complete solution to the
K\"ahler-Ricci flow on a complex  manifold $M$ with bounded
curvature and nonnegtive bisectional curvature and $0\leq t < T$.
For any point $x\in M$ and any vector $V$ in the holomorphic
tangent space $T^{1,0}_xM$, let
$$
Q_{\alpha\bar \beta}=\frac{\partial}{\partial t}R_{\alpha\bar
\beta}+ R_{\alpha\bar \gamma}R_{\gamma\bar \beta}
+\nabla_{\gamma}R_{\alpha\bar \beta} V^{\gamma}
+\nabla_{\bar\gamma}R_{\alpha\bar
\beta}V^{\bar\gamma}+R_{\alpha\bar \beta\gamma\bar \delta}V^{
\gamma}V^{\bar\delta}+\frac {1}{t}R_{\alpha\bar \beta} .
$$
Then we have
$$
Q_{\alpha\bar \beta}W^{\alpha}W^{\bar \beta}\geq 0
$$
for all $x\in M$, $V, W\in T^{1,0}_xM$, and $t>0$.
\end{theorem}

\begin{corollary}[Kähler–Ricci scalar and determinant estimates]
  \label{cor:cz-kahlerricci-scalar-and-determinant-estimates}
  \source{cao-zhu}
  \dcref{Corollary 2.5.10}
  \group{Cao--Zhu Li-Yau-Hamilton Estimates}
Under the assumptions of Theorem $2.5.9,$ we have
\begin{itemize}
\item[(i)] the scalar curvature $R$ satisfies the estimate
$$
\frac{\partial R}{\partial t}-\frac{|\nabla R|^2}{R}
+\frac{R}{t}\geq 0,$$ and \item[(ii)] assuming $R_{\alpha\bar
\beta}>0$, the determinant
$\phi=\det(R_{\alpha\bar\beta})/\det(g_{\alpha\bar\beta})$ of the
Ricci curvature satisfies the estimate
$$
\frac {\partial \phi}{\partial t}-\frac{|\nabla\phi|^2}{n\phi}
+\frac {n\phi}{t}\geq 0
$$
for all $x\in M$ and $t>0$.
\end{itemize}
\end{corollary}

\section{Perelman's Estimate for Conjugate Heat Equations}

Write $\tau=T-t$.  The conjugate heat equation used below is
\[
\partial_\tau u=\Delta u-Ru,
\tag{2.6.7}
\]
with $u=(4\pi\tau)^{-n/2}e^{-f}$.

\begin{lemma}[gradient estimate for conjugate heat equations]
  \label{lem:cz-gradient-estimate-for-conjugate-heat-equations}
  \source{cao-zhu}
  \dcref{Lemma 2.6.1}
  \group{Cao--Zhu Perelman's Estimate for Conjugate Heat Equations}
Let $g_{ij}(x,t)$, $0 \leq t < T$, be a complete solution to the
Ricci flow on an $n$-dimensional manifold $M$ and let $u=(4 \pi
\tau)^{-\frac{n}{2}} e^{-f}$ be a solution to the conjugate
equation $(2.6.7)$ with $\tau=T-t$. Set
$$
H=2 \Delta f -| \nabla f |^2 + R + \frac{f-n}{\tau}
$$
and
$$
v = \tau H u = \big(\tau(R+2 \Delta f - | \nabla f|^2)+f-n \big)
u.
$$
Then we have
$$
\frac{\partial H}{\partial \tau} = \Delta H - 2 \nabla f \cdot
\nabla H - \frac{1}{\tau} H - 2 \left| R_{ij} + \nabla_i\nabla_j f
- \frac{1}{2 \tau} g_{ij} \right|^2,
$$
and
$$
\frac{\partial v}{\partial \tau} = \Delta v - Rv - 2 \tau u
\left|R_{ij}+\nabla_i\nabla_jf - \frac{1}{2 \tau} g_{ij}\right|^2.
$$
\end{lemma}

\begin{proof}
  \sketch Differentiate $H$ and $v=\tau Hu$, use the conjugate heat equation and the Ricci-flow evolution identities, and complete the square.
\end{proof}
